% R5 — Training strategy
\subsection{Training strategy}
The encoder is trained on the Normal-, MS-, and PD-labeled videos in the training partition, without using their labels in the feature-learning objective. Each batch samples source videos with equal probability and then samples windows within those sources. This prevents recordings with more windows from dominating encoder training. Batches contain 32 windows.

The context encoder receives an augmented view with parts of the skeleton hidden, and the predictor estimates the target encoder's features at the hidden positions. Training compares distributions over feature dimensions using cross-entropy with centering and sharpening~\cite{dino}. These distributions describe learned features, rather than the three collection labels. The target encoder is updated gradually from the context encoder's weights. Spatial augmentation uses rotation within $\pm15^\circ$, scaling from 0.9 to 1.1, translation within $\pm0.1$ normalized units, and horizontal reflection with probability 0.5. A single transformation is applied throughout each window. Reflection exchanges left and right landmark slots; its effect on useful gait asymmetry was not separately evaluated.

We train for 800 updates and evaluate a further 400-update continuation from the learned weights using the same training sources and objective. The continuation restarts the optimizer, running center, random-number generators and schedules. It therefore differs from uninterrupted training for 1,200 updates. For each stage, a linear classifier is fitted to the frozen features. The stage with higher source-weighted validation macro-F1 is selected, retaining the initial stage on a tie. Validation videos are excluded from encoder and classifier fitting, and test predictions are generated only after stage selection. The selected model is not refitted using the validation videos.
